%! TeX root = ../main.tex
\chapter{PENDAHULUAN}
\label{chap:pendahuluan}

\section{Latar Belakang}
\label{sec:latarbelakang}

\lipsum[1-2]

\section{Rumusan Masalah}
\label{sec:rumusanmasalah}

Berdasarkan latar belakang yang telah diuraikan, rumusan masalah dalam penelitian ini adalah:
\begin{enumerate}
  \item Bagaimana \lipsum[1][1-2]
  \item Bagaimana \lipsum[2][1-2]
\end{enumerate}

\section{Batasan Masalah}
\label{sec:batasanmasalah}

Batasan masalah dalam penelitian ini adalah sebagai berikut:
\begin{enumerate}
  \item \lipsum[3][1-2]
  \item \lipsum[4][1-2]
\end{enumerate}

\section{Tujuan}
\label{sec:tujuan}

Tujuan yang ingin dicapai dari penelitian ini adalah:
\begin{enumerate}
  \item \lipsum[1][1-2]
  \item \lipsum[2][1-2]
\end{enumerate}

\section{Manfaat}
\label{sec:manfaat}

Manfaat yang diharapkan dari hasil penelitian ini antara lain:
\begin{enumerate}
  \item \textbf{Bagi Peneliti dan Akademisi}: \lipsum[1][1-3]
  \item \textbf{Bagi Sektor Kesehatan/Klinis}: \lipsum[2][1-3]
\end{enumerate}

